% !TEX root = task-2-10-6.tex
% Editable task bank. Append new BankTask blocks at the end of the document.
\documentclass[12pt,a4paper]{article}
\usepackage{physics-textbook}
\usepackage{graphicx,needspace}
\usetikzlibrary{decorations.pathmorphing}
\newif\ifShowSolutions
\ifdefined\TutorEdition
  \ShowSolutionstrue
\else
  \ShowSolutionsfalse % Change to \ShowSolutionstrue to display solutions.
\fi
\newcounter{banktask}
\newsavebox{\BankTaskBox}
\newcommand{\BankTaskContent}[3]{%
  \noindent\textbf{\thebanktask.}\quad #2\par
  \ifShowSolutions
    \Needspace{50mm}% Keep the solution heading with its opening text or diagram.
    #3\par
    \ifdefined\TutorEdition
      \ifcsname TutorSource#1\endcsname
        \nopagebreak\smallskip{\small\csname TutorSource#1\endcsname\par}%
      \fi
    \fi
  \fi
}
\newcommand{\BankTask}[3]{%
  \refstepcounter{banktask}\par\medskip
  \sbox{\BankTaskBox}{\begin{minipage}{\linewidth}%
    \BankTaskContent{#1}{#2}{#3}\end{minipage}}%
  \ifdim\dimexpr\ht\BankTaskBox+\dp\BankTaskBox\relax>\dimexpr\textheight-20mm\relax
    \Needspace{35mm}\BankTaskContent{#1}{#2}{#3}%
  \else
    \Needspace{\dimexpr\ht\BankTaskBox+\dp\BankTaskBox\relax}%
    \noindent\usebox{\BankTaskBox}\par
  \fi
  \medskip
}
\begin{document}
\BookSubsection{2.10.6}{Кипение}


% Source 055098: https://ege.fipi.ru/bank/questions.php?search=1&proj=BA1F39653304A5B041B656915DC36B38&qid=055098
\BankTask{055098}{
Стальную кастрюлю, наполовину заполненную водой, поставили на плиту, закрыли лёгкой крышкой и стали нагревать воду до кипения. Атмосферное давление составляет \(760\,\text{мм рт. ст.}\)
\begin{enumerate}
\item Относительная влажность воздуха под крышкой в процессе нагревания воды уменьшается.
\item Внутренняя энергия водяного пара в кастрюле в процессе нагревания воды увеличивается.
\item При нагревании воды выше \(80\,^{\circ}\mathrm{C}\), когда схлопываются образующиеся пузырьки, давление водяных паров постепенно уменьшается.
\item Температура дна кастрюли вблизи границы с водой при кипении немного превышает \(100\,^{\circ}\mathrm{C}\).
\item Плотность насыщенных водяных паров над поверхностью воды в процессе её нагревания до температуры кипения увеличивается.
\end{enumerate}
}{
\smallskip\noindent\textbf{Решение от GPT}\par
1) Неверно: под крышкой воздух насыщается водяным паром, его относительная влажность не уменьшается.

2) Верно: при нагревании растут температура и количество водяного пара, поэтому его внутренняя энергия увеличивается.

3) Неверно: давление насыщенного пара с повышением температуры увеличивается; пузырьки схлопываются при попадании в более холодные слои воды вследствие конденсации пара.

4) Верно: при данном атмосферном давлении вода кипит при \(100\,^{\circ}\mathrm{C}\); для передачи теплоты от дна воде температура дна должна быть немного выше.

5) Верно: плотность насыщенного водяного пара увеличивается с температурой.

Ответ: \(245\).
}

% Source 25410: https://phys-ege.sdamgia.ru/problem?id=25410
\BankTask{25410}{
В вертикальном цилиндре, закрытом легким поршнем, находится бензол ($\mathrm{C_6H_6}$) при температуре кипения t  =  $80$ \ensuremath{{}^\circ}C. При сообщении бензолу количества теплоты Q часть его превращается в пар, который при изобарном расширении совершает работу А. Удельная теплота парообразования бензола L  =  $396$ · $10^{3}$ Дж/кг, его молярная масса M  =  $78$ · $10^{-3}$ кг/моль. Какая часть подведенного к бензолу количества теплоты переходит в работу? Объемом жидкого бензола пренебречь.
}{
\smallskip\noindent\textbf{Решение}\par
$1$.   При кипении бензола происходит его изобарное расширение. Работа пара $A = p\Delta V,$ где p  — атмосферное давление, $\Delta V$  — изменение объема.
$2$.  Считая пар идеальным газом, воспользуемся уравнением Клапейрона  — Менделеева для определения изменения объема за счет испарившегося бензола массой $\Delta m:$ $p\Delta V= \frac{\Delta m}{M} RT,$ где M  =  $78$ · $10^{-3}$ кг/моль  — молярная масса бензола, T  =  $80$ + $273$  =  $353$ K  — температура кипения бензола. Отсюда $A= \frac{\Delta mRT}{M} .$
$3$.  Количество теплоты Q, необходимое для испарения массы $\Delta m$ бензола, пропорционально удельной теплоте парообразования L: $Q = \Delta mL .$
$4$.  Искомая величина определяется отношением:

\[\eta = \frac{A}{Q} = \frac{RT}{ML} = \frac{8{,}31 \cdot 353}{78 \cdot 10 ^{ - 3 } \cdot 396 \cdot 10 ^3} \approx 0{,}095.\]

Ответ:$\eta \approx 0{,}095.$
}

\end{document}
